\documentclass[
	fontsize=12pt,          % default font size 12pt
	paper=a4,               % DIN A4 page format
	numbers=noenddot,       % remove dots behind chapter numbers (e.g. 1.5 not 1.5.)
	listof=totoc,           % add list of figures, tables, etc. to ToC
	listof=entryprefix,     % add entry name to figures, tables, etc.
	listof=nochaptergap,    % no chapter gap for figures, tables, etc.
	bibliography=totoc,     % add bibliography to ToC but without a chapter number
	parskip=half            % half line spacing between paragraphs
	openany                 % chapters can start on any page
]{scrbook}
\usepackage{tcolorbox}
\usepackage{tabularx}

\include{../../../for_latex/preamble}
\title{\Huge\textbf{Längster Stab (Tutor)}}
\author{Luis Staudt}
\date{}

% Code-Highlighting für Java
\definecolor{codegray}{rgb}{0.5,0.5,0.5}
\definecolor{codepurple}{rgb}{0.58,0,0.82}
\definecolor{backcolour}{rgb}{0.95,0.95,0.92}
\definecolor{keywordcolor}{rgb}{0.8,0.47,0.13}

\lstdefinestyle{java}{
	language=,
	basicstyle=\ttfamily\small,
	keywordstyle=\bfseries\color{blue!70!black},
	stringstyle=\color{red!60!black},
	commentstyle=\itshape\color{gray},
	numbers=none,
	numberstyle=\tiny\color{gray},
	numbersep=8pt,
	frame=single,
	framerule=0.4pt,
	rulecolor=\color{gray!50},
	breaklines=true,
	showstringspaces=false,
	tabsize=4,
	xleftmargin=1.5em,
	framexleftmargin=1.5em,
	aboveskip=0.8em,
	belowskip=0.6em
}
\lstset{style=java}

\newtcolorbox{hinweis}[1][Hinweis]{
	colback=blue!5,
	colframe=blue!40!black,
	fonttitle=\bfseries,
	title={#1},
	boxrule=0.5pt,
	arc=2pt
}

\newtcolorbox{beispiel}[1][Beispiel]{
	colback=green!5,
	colframe=green!40!black,
	fonttitle=\bfseries,
	title={#1},
	boxrule=0.5pt,
	arc=2pt
}

\newtcolorbox{loesung}[1][Musterlösung]{
	colback=orange!5,
	colframe=orange!55!black,
	fonttitle=\bfseries,
	title={#1},
	boxrule=0.5pt,
	arc=2pt
}

\begin{document}

% --- Titelblock ---
\begin{center}
	\rule{\textwidth}{0.8pt}\\[0.6em]
	{\LARGE\bfseries Intensivtutorium}\\[0.3em]
	{\large Programmierung \quad\textbar\quad Tutorexemplar mit Musterlösung}\\[0.6em]
	\begin{tabular}{r l @{\hspace{2cm}} r l}
		\textbf{Semester:}      & Sommersemester 2026                    & \textbf{Tutor:}         & Luis Staudt \\
		\textbf{Studiengänge:}  & PEB \& MVB         & & \\
	\end{tabular}\\[0.6em]
	\rule{\textwidth}{0.8pt}
\end{center}
\vspace{1em}

% --- Seitenkopf und -fuß (ab Seite 2) ---
\pagestyle{fancy}
\fancyhf{}
\lhead{\small Intensivtutorium, Tutorexemplar}
\rhead{\small SS 2026}
\cfoot{\thepage}
\renewcommand{\headrulewidth}{0.4pt}


% =============================================
% Aufgabe 4: Längster Stab (Tutor)
% =============================================

\section*{Aufgabe 4: Längster Stab \;\normalsize(Tutorexemplar)}

\begin{hinweis}[Lernziel der Aufgabe]
	Diese Aufgabe ist im Stil der \emph{Zusatzaufgaben} zum Fachwerk-Praktikum gestaltet. Der eigentliche Lernkern ist nicht das Schreiben von neuem Code, sondern:
	\begin{itemize}[nosep]
		\item \textbf{Lesen und Debuggen fremden, fehlerhaften Codes},
		\item \textbf{Nachschlagen in einer API-Dokumentation} (\texttt{fachwerk\_API.md}),
		\item das saubere Umsetzen einer \textbf{Maximum-Suche} kombiniert mit dem \textbf{Satz des Pythagoras}.
	\end{itemize}
\end{hinweis}

\subsection*{Aufgabenstellung (Kurzfassung)}

Für das vom Anwender gewählte Fachwerk (0 bis 4) soll die Stabnummer des \textbf{längsten Stabes} und dessen Länge ausgegeben werden. Die Stablänge folgt aus den Koordinaten von Anfangs- und Endknoten:
\[
	l = \sqrt{(x_e - x_a)^2 + (y_e - y_a)^2}
\]
Ein Aufruf von \texttt{berechneGroessen()} ist \textbf{nicht} nötig (reine Geometrie).

\subsection*{Eingebaute Fehler im Vorlagencode}

Der ausgegebene Vorlagencode enthält \textbf{acht} bewusst platzierte Fehler über die typischen Fehlerklassen der Zusatzaufgaben:

\begin{center}
\small
\renewcommand{\arraystretch}{1.3}
\begin{tabularx}{\textwidth}{@{}c >{\raggedright\arraybackslash}X >{\raggedright\arraybackslash}X >{\raggedright\arraybackslash}X@{}}
	\toprule
	\textbf{\#} & \textbf{Fehler im Vorlagencode} & \textbf{Korrektur} & \textbf{Fehlerklasse} \\
	\midrule
	1 & \texttt{gibAnzahllStaebe()} & \texttt{gibAnzahlStaebe()} & Tippfehler (Methodenname) \\
	2 & \texttt{double a}, \texttt{double e} & \texttt{int a}, \texttt{int e} & Falscher Typ (Index als \texttt{double}) \\
	3 & \texttt{dy = \dots gibX() \dots gibY()} & beide \texttt{gibY()} & Falsche Methode (X statt Y) \\
	4 & \texttt{dx*dx + dy + dy} & \texttt{dx*dx + dy*dy} & Logikfehler (\texttt{+} statt \texttt{*}) \\
	5 & fehlendes \texttt{;} nach \texttt{Math.sqrt(\dots)} & Semikolon ergänzen & Syntaxfehler \\
	6 & fehlendes \texttt{;} nach \texttt{nr = i} & Semikolon ergänzen & Syntaxfehler \\
	7 & \texttt{System.out.pritnln} & \texttt{println} & Tippfehler \\
	8 & \texttt{Nr} in der Ausgabe & \texttt{nr} & Falsche Variable (Groß/Klein) \\
	\bottomrule
\end{tabularx}
\end{center}

\begin{hinweis}[Häufige Stolpersteine bei der Besprechung]
	\begin{itemize}[nosep]
		\item Fehler~2 fällt oft erst beim Kompilieren auf: \texttt{gibKnoten(int)} akzeptiert keinen \texttt{double}. Hier lohnt der Blick in die API.
		\item Fehler~3 und~4 sind \emph{logische} Fehler. Das Programm kompiliert (nach Behebung der Syntaxfehler) und läuft, liefert aber falsche Längen. Gut geeignet, um auf das \emph{Testen} mit bekannten Werten hinzuweisen.
		\item Es genügt, $l^2$ zu vergleichen; das Ziehen der Wurzel ist für die reine Maximum-Suche nicht nötig. Für die \emph{Ausgabe} der Länge wird sie aber gebraucht.
	\end{itemize}
\end{hinweis}

\subsection*{Musterlösung}

\begin{loesung}[Korrigierter Quelltext]
\begin{lstlisting}
public class Fachwerkaufgabe {
    public static void main(String[] args) {
        Fachwerk fw = null;
        System.out.println("Welches Fachwerk wollen Sie rechnen (0-4)?");
        int wahl = Tastatureingabe.leseInt();
        fw = new Fachwerk(wahl);

        double max = 0;
        int nr = 0;

        for (int i = 0; i < fw.gibAnzahlStaebe(); i++) {
            int a = fw.gibStab(i).gibKnotenNrA();
            int e = fw.gibStab(i).gibKnotenNrE();
            double dx = fw.gibKnoten(e).gibX() - fw.gibKnoten(a).gibX();
            double dy = fw.gibKnoten(e).gibY() - fw.gibKnoten(a).gibY();
            double laenge = Math.sqrt(dx * dx + dy * dy);
            if (laenge > max) {
                max = laenge;
                nr = i;
            }
        }
        System.out.println("Laengster Stab: Stab " + nr + " mit Laenge " + max);
    }
}
\end{lstlisting}
\end{loesung}

Eine lauffähige Variante im Package \texttt{prog} liegt unter \texttt{src/prog/LongestBar.java}.

\begin{beispiel}[Mögliche Ausgabe (Fachwerk~2)]
\begin{verbatim}
Welches Fachwerk wollen Sie rechnen (0-4)?
2
Laengster Stab: Stab 7 mit Laenge 4.272001872...
\end{verbatim}
\emph{Hinweis:} Die konkrete Stabnummer und Länge hängen vom gewählten Fachwerk
(0 bis 4) ab und sollten beim Korrigieren am tatsächlichen Fachwerk geprüft werden.
\end{beispiel}

\subsection*{Mögliche Erweiterungen (für schnelle Gruppen)}

\begin{itemize}[nosep]
	\item Zusätzlich den \textbf{kürzesten} Stab ausgeben (Initialisierung von \texttt{min} mit dem ersten Stab oder \texttt{Double.MAX\_VALUE}).
	\item Die \textbf{Gesamtlänge} aller Stäbe (Materialbedarf) aufsummieren.
	\item Den \textbf{Mittelpunkt} des längsten Stabes mit ausgeben (Brücke zu Zusatzaufgabe~3).
\end{itemize}

\end{document}
